\SourceLeaf{302}
obtinetur, ita ut forma negativa $(-1, 0, D)$ semper habeat characterem ex $Q$.
Sed quoniam haec forma cum quacunque alia pr.\ primitiva negativa eiusdem det.
composita similem formam positivam producit, facile perspicitur, nullam formam
pr.\ prim.\ negativam characterem ex $P$ habere posse.

II. Pro generibus improprie primitivis (positivis) simili modo probatur, rem
vel eodem modo se habere ut in proprie primitivis, vel contrario, prout $D\equiv1$
vel $\equiv5\;(\mathrm{mod.}\,8)$. Nam in casu priori erit etiam $D'\equiv1\;(\mathrm{mod.}\,8)$, unde facile
concluditur, inter numeros $a$, $b$, $c$ etc.\ vel nullum numerum formae $8n+3$ et
$8n+5$ reperiri vel duos vel quatuor etc.\ (scilicet productum ex quotcunque
numeris imparibus, inter quos numeri formae $8n+3$ et $8n+5$ coniunctim
multitudinem imparem efficiunt, semper evadit vel $\equiv3$ vel $\equiv5\;(\mathrm{mod.}\,8)$, productum
autem ex omnibus $a$, $b$, $c$ etc.\ aequale esse debet vel ipsi $D'$ vel ipsi $-D'$); hinc
patet, characterem integrum formae $(2, 1, \frac{1-D}{2})$ involvere vel nullum characterem
particularem ex $\Omega$, vel duos vel quatuor etc., adeoque pertinere ad $P$. Iam
quum quaevis forma improprie primitiva (positiva) determinantis $D$ spectari possit
tamquam composita ex $(2, 1, \frac{1-D}{2})$ atque proprie primitiva (positiva) eiusdem
determinantis, perspicuum est, nullam formam improprie primitivam (positivam)
characterem ex $Q$ in hoc casu habere posse. In casu altero. $D\equiv5\;(\mathrm{mod.}\,8)$,
omnia contraria sunt, scilicet $D'$, qui etiam erit $\equiv5$, certo multitudinem
imparem factorum formae $8n+3$ atque $8n+5$ implicabit, unde concluditur,
characterem formae $(2, 1, \frac{1-D}{2})$, atque hinc etiam characterem cuiusvis formae
improprie primitivae (pos.) det.\ $D$ pertinere ad $Q$, adeoque nulli characterum
$P$ genus impr.\ prim.\ pos.\ respondere posse.

III. Denique pro determinante negativo genera improprie primitiva
negativa rursus contraria sunt generibus improprie primitivis positivis, scilicet illa non
poterunt habere characterem ex $P$ vel ex $Q$, prout $D\equiv1$ vel $\equiv5\;(\mathrm{mod.}\,8)$,
sive prout $-D$ est formae $8n+7$ vel $8n+3$. Hoc nullo negotio deducitur
inde, quod ex compositione formae $(-1, 0, D)$, cuius character est ex $Q$, cum
formis improprie primitivis negativis eiusdem determinantis formae improprie
primitivae positivae proveniunt, adeoque, quando ab his exclusi sunt characteres $Q$,
necessario ab illis exclusi esse debent characteres $P$, et contra.

\Needspace{11\baselineskip}
\SourceLeaf{303}
\begin{center}
\textit{Methodus peculiaris, numeros primos in duo quadrata decomponendi.}
\end{center}

\Needspace{8\baselineskip}
\Article{265}
Ex disquisitionibus artt.\ 257, 258 supra multitudine classium ancipitum,
quibus omnia praecedentia sunt superstructa, multae aliae conclusiones attentione
perdignae deduci possunt, quas brevitatis caussa supprimere oportet; sequentem
tamen, elegantia sua insignem, praeterire non possumus. Pro determinante
positivo $p$, qui est numerus primus formae $4n+1$, unicam tantummodo
classem ancipitem proprie primitivam dari ostendimus; quapropter omnes formae
ancipites proprie primitivae talis determinantis proprie aequivalentes erunt. Si
itaque $b$ est numerus integer positivus proxime minor quam $\surd p$, atque $p-bb
=a'$, formae $(1, b, -a')$, $(-1, b, a')$, proprie aequivalebunt, adeoque, quum utraque
manifesto sit forma reducta, altera in alterius periodo erit contenta. Tribuendo
formae priori in periodo sua indicem $0$, index posterioris necessario erit impar (quoniam
termini primi harum duarum formarum signa opposita habent); ponatur itaque
$=2m+1$. Porro facile perspicitur, si formae indicum $1, 2, 3$ etc.\ resp.\ sint
\[
(-a', b', a''),\quad (a'', b'', -a'''),\quad (-a''', b''', a'''')\text{ etc.}
\]
indicibus $2m$, $2m-1$, $2m-2$, $2m-3$ etc.\ responsuras esse formas
\[
(a', b, -1),\quad (-a'', b', a'),\quad (a''', b'', -a''),\quad (-a'''', b''', a''')\text{ etc.}
\]
Hinc colligitur, si forma indicis $m$ sit $(A, B, C)$, eandem fore $(-C, B, -A)$,
adeoque $C=-A$ et $p=BB+AA$. Quare quivis numerus primus formae
$4n+1$ in duo quadrata decomponi potest (quam propositionem supra, art.\ 182,
e principiis prorsus diversis deduximus), et ad talem decompositionem pervenire
possumus per methodum simplicissimam et omnino uniformem, scilicet per evolutionem
periodi formae reductae, cuius determinans est ille numerus primus et cuius
terminus primus $1$, usque ad formam, cuius termini externi magnitudine sunt
aequales, signis oppositi. Ita \textit{e.g.}\ pro $p=233$ habetur $(1, 15, -8)$, $(-8, 9, 19)$,
$(19, 10, -7)$, $(-7, 11, 16)$, $(16, 5, -13)$, $(-13, 8, 13)$, atque $233=64+169$.
Ceterum patet, $A$ necessario fieri imparem (quoniam $(A, B, -A)$ debet esse
forma proprie primitiva), et proin $B$ parem. --- Quum pro determinante positivo
$p$, qui est numerus primus formae $4n+1$, etiam in ordine improprie primitivo
unica tantum classis anceps contineatur, perspicuum est, si $g$ sit numerus
\SourceLeaf{304}
impar proxime minor quam $\surd p$, atque $p-gg=4h$, formas reductas improprie
primitivas $(2, g, -2h)$, $(-2, g, 2h)$ proprie aequivalere, adeoque alteram
in alterius periodo contentam esse. Hinc per ratiocinia praecedentibus omnino
similia concluditur, in periodo formae $(2, g, -2h)$ reperiri formam, cuius termini
externi magnitudine aequales sint, signa habeant opposita, ita ut discerptio
numeri $p$ in duo quadrata etiam hinc peti possit. Patet autem, terminos externos
huius formae fore pares, adeoque medium imparem; et quum constet, numerum
primum unico tantum modo in duo quadrata decomponi posse, forma per
hanc posteriorem methodum inventa erit vel $(B, \pm A, -B)$ vel $(-B, \pm A, B)$.
Ita in exemplo nostro pro $p=233$ habetur $(2, 15, -4)$, $(-4, 13, 16)$, $(16, 3, -14)$,
$(-14, 11, 8)$, $(8, 13, -8)$, et $233=169+64$ ut supra.

\Needspace{11\baselineskip}
\begin{center}
\textit{DIGRESSIO CONTINENS TRACTATUM DE FORMIS TERNARIIS.}
\end{center}

\Needspace{8\baselineskip}
\Article{266}
Hactenus disquisitionem nostram ad tales functiones secundi gradus restrinximus,
quae \textit{duas} indeterminatas implicant, neque opus fuit, denominationem
specialem ipsis tribuere. Sed manifesto hoc argumentum tamquam sectionem
maxime particularem disquisitionis generalissimae \textit{de functionibus algebraicis
rationalibus integris h\,o\,m\,o\,g\,e\,n\,e\,i\,s plurium indeterminatarum et plurium dimensionum}
considerare, talesque functiones secundum multitudinem dimensionum in \textit{formas
secundi, tertii, quarti gradus etc.}, secundum multitudinem indeterminatarum autem
in \textit{formas binarias, ternarias, quaternarias} etc.\ commode distinguere possumus.
Formae itaque, hactenus simpliciter sic dictae, vocabuntur \textit{formae binariae
secundi gradus}; tales autem functiones ut
\[
Axx+2Bxy+Cyy+2Dxz+2Eyz+Fzz
\]
denotantibus $A$, $B$, $C$, $D$, $E$, $F$ integros datos) dicentur \textit{formae ternariae secundi
gradus} et sic porro. Proxime quidem Sectio praesens solis formis binariis secundi
gradus est dicata; sed quoniam complures veritates ad has spectantes, eaeque pulcherrimae,
adhuc supersunt, quarum fons proprius in theoria formarum ternariarum
secundi gradus est quaerendus, brevem ad hanc theoriam digressionem hic intercalamus,
in qua ex primis eius elementis ea trademus, quae ad perfectionem
theoriae formarum binariarum sunt necessaria, quod geometris acceptius fore spe%
\SourceLeaf{305}%
ramus, quam si illas vel supprimeremus, vel per methodos minus genuinas erueremus.
Exactiorem autem de hoc argumento gravissimo disquisitionem ad aliam
occasionem nobis reservare debemus, tum quod ipsius ubertas limites huius operis
iam nunc longe egrederetur, tum quod spes est, luculentis adhuc incrementis
eam in posterum locupletatum iri. Formae vero tum quaternariae, quinariae etc.
secundi gradus, tum omnes superiorum graduum hoc quidem loco ab instituto
nostro penitus excluduntur\HistoricalNote{*)}{Propter hanc rationem formae binariae vel ternariae \textit{secundi gradus} in sequentibus semper sunt intelligendae, quoties de talibus formis simpliciter loquemur.}, sufficiatque hunc campum vastissimum geometrarum
attentioni commendavisse, in quo materiem ingentem vires suas exercendi,
Arithmeticamque sublimiorem egregiis incrementis augendi inveni\`ent.

\Needspace{8\baselineskip}
\Article{267}
Ad perspicuitatem multum proderit, inter tres indeterminatas, in formam
ternariam ingredientes, simili modo ut in formis binariis, ordinem fixum stabilire,
ita ut \textit{indeterminata prima, secunda} et \textit{tertia} ab invicem distinguantur; in disponendis
autem singulis formae partibus hunc ordinem semper observabimus, ut
primum locum obtineat ea pars quae quadratum indeterminatae primae implicat,
in sequentibus eae quae implicant quadratum indeterminatae secundae, quadratum
tertiae, productum duplum secundae in tertiam, productum duplum primae
in tertiam, productum duplum primae in secundam deinceps sequantur; denique
numeros integros determinatos per quos haec quadrata et producta dupla multiplicata
sunt, eodem ordine \textit{co\"efficientem primum, secundum, tertium, quartum, quintum,
sextum} vocabimus. Ita
\[
axx+a'x'x'+a''x''x''+2bx'x''+2b'xx''+2b''xx'
\]
erit forma ternaria rite ordinata, cuius indeterminata prima $x$, secunda $x'$, tertia
$x''$, co\"efficiens primus $a$ etc., quartus $b$ etc.\ Sed quoniam ad brevitatem multum
conferet, si non semper necesse est, indeterminatas formae ternariae per literas
peculiares denotare, eandem formam, quatenus ad indeterminatas non respicimus,
etiam hoc modo
\[
\begin{pmatrix}
a, & a', & a'' \\
b, & b', & b''
\end{pmatrix}
\]
designabimus.
\SourceLeaf{306}

Ponendo
\[
\begin{array}{lll}
bb-a'a''=A, & b'b'-aa''=A', & b''b''-aa'=A'' \\
ab-b'b''=B, & a'b'-bb''=B', & a''b''-bb'=B''
\end{array}
\]
oritur alia forma
\[
\begin{pmatrix}
A, & A', & A'' \\
B, & B', & B''
\end{pmatrix}\ldots F
\]
quam \textit{formae}
\[
\begin{pmatrix}
a, & a', & a'' \\
b, & b', & b''
\end{pmatrix}\ldots f
\]
\textit{adiunctam} dicemus. Hinc rursus invenitur, denotando brevitatis caussa numerum
\[
abb+a'b'b'+a''b''b''-aa'a''-2bb'b''\quad\text{per}\quad D
\]
\[
\begin{array}{lll}
BB-A'A''=aD, & B'B'-AA''=a'D, & B''B''-AA'=a''D \\
AB-B'B''=bD, & A'B'-BB''=b'D, & A''B''-BB'=b''D
\end{array}
\]
unde patet, formae $F$ adiunctam esse formam
\[
\begin{pmatrix}
aD, & a'D, & a''D \\
bD, & b'D, & b''D
\end{pmatrix}
\]
Numerum $D$, a cuius indole proprietates formae ternariae $f$ imprimis pendent,
\textit{determinantem} huius formae vocabimus; hoc modo determinans formae $F$ fit
$=DD$, sive aequalis quadrato determinantis formae $f$, cui adiuncta est.

Ita \textit{e.\ g.}\ formae ternariae $\left(\begin{smallmatrix}29,\allowbreak &13,\allowbreak &9\\7,\allowbreak &-1,\allowbreak &14\end{smallmatrix}\right)$
adiuncta est $\left(\begin{smallmatrix}-68,\allowbreak &-260,\allowbreak &-181\\217,\allowbreak &-111,\allowbreak &133\end{smallmatrix}\right)$,
utriusque determinans $=1$.

Formae ternariae determinantis $0$ ab investigatione sequente omnino excludentur,
quippe quae, ut in formarum ternariarum theoria, alia occasione uberius
tradenda, ostendetur, \textit{specie} tantum sunt ternariae, reveraque binariis aequipollentes.

\Needspace{8\baselineskip}
\Article{268}
Si forma aliqua ternaria $f$ determinantis $D$, cuius indeterminatae sunt
$x$, $x'$, $x''$ (puta prima $=x$ etc.) in formam ternariam $g$ determinantis $E$, cuius
indeterminatae sunt $y$, $y'$, $y''$, transmutatur per substitutionem talem
\SourceLeaf{307}
\[
\begin{aligned}
x&=\alpha y+\beta y'+\gamma y'' \\
x'&=\alpha'y+\beta'y'+\gamma'y'' \\
x''&=\alpha''y+\beta''y'+\gamma''y''
\end{aligned}
\]
ubi novem co\"efficientes $\alpha$, $\beta$ etc.\ omnes supponuntur esse numeri integri, brevitatis
caussa neglectis indeterminatis simpliciter dicemus, $f$ transire in $g$ per substitutionem
$(S)$
\[
\begin{array}{ccc}
\alpha, & \beta, & \gamma \\
\alpha', & \beta', & \gamma' \\
\alpha'', & \beta'', & \gamma''
\end{array}
\]
atque $f$ \textit{implicare} ipsam $g$, sive $g$ \textit{sub $f$ contentam} esse. Ex tali itaque suppositione
sponte sequuntur sex aequationes pro sex co\"efficientibus in $g$, quas apponere
non erit necessarium; hinc autem per calculum facilem sequentes conclusiones
evolvuntur:

I. Designato brevitatis caussa numero
\[
\alpha\beta'\gamma''+\beta\gamma'\alpha''+\gamma\alpha'\beta''
-\gamma\beta'\alpha''-\alpha\gamma'\beta''-\beta\alpha'\gamma''\quad\text{per }k
\]
invenitur post debitas reductiones $E=kkD$, unde patet, $D$ metiri ipsum $E$
et quotientem esse quadratum. Patet itaque, numerum $k$ pro transformationibus
formarum ternariarum simile quid esse, ac numerum $\alpha\delta-\beta\gamma$ in art.\ 157
pro transformationibus formarum binariarum, puta radicem quadratam ex quotiente
determinantium, unde coniectare possemus, diversitatem signi ipsius $k$
etiam hic stabilire differentiam essentialem inter transformationes atque implicationes
proprias et improprias. Sed rem propius contemplando perspicuum est. $f$
transire in $g$ etiam per hanc substitutionem
\[
\begin{array}{ccc}
-\alpha, & -\beta, & -\gamma \\
-\alpha', & -\beta', & -\gamma' \\
-\alpha'', & -\beta'', & -\gamma''
\end{array}
\]
ponendo autem in valore ipsius $k$ pro $\alpha$, $-\alpha$, pro $\beta$, $-\beta$ etc.\ prodibit $-k$,
quare haec substitutio substitutioni $S$ dissimilis foret, et quaevis forma ternaria,
aliam uno modo implicans, eandem etiam altero modo implicaret. Talis itaque
\SourceLeaf{308}
distinctio, quoniam in formis ternariis nullum usum habet, hic omnino proscribetur.

II. Denotando per $F$, $G$ formas ipsis $f$, $g$ resp.\ adiunctas, determinantur
co\"efficientes in $F$ per co\"efficientes in $f$, co\"efficientesque in $G$ per valores
co\"efficientium formae $g$ ex aequationibus quas suppeditat substitutio $S$ notos.
Exprimendo co\"efficientes formae $f$ per literas, ex comparatione valorum co\"efficientium
formarum $F$, $G$ nullo negotio confirmatur, $F$ implicare formam $G$ atque
in eam transmutari per substitutionem $(S')$
\[
\begin{array}{ccc}
\beta'\gamma''-\beta''\gamma', & \gamma'\alpha''-\gamma''\alpha', & \alpha'\beta''-\alpha''\beta' \\
\beta''\gamma-\beta\gamma'', & \gamma''\alpha-\gamma\alpha'', & \alpha''\beta-\alpha\beta'' \\
\beta\gamma'-\beta'\gamma, & \gamma\alpha'-\gamma'\alpha, & \alpha\beta'-\alpha'\beta
\end{array}
\]
Calculum ipsum nullis difficultatibus obnoxium non adscribimus.

III. Forma $g$ per substitutionem $(S'')$
\[
\begin{array}{ccc}
\beta'\gamma''-\beta''\gamma', & \beta''\gamma-\beta\gamma'', & \beta\gamma'-\beta'\gamma \\
\gamma'\alpha''-\gamma''\alpha', & \gamma''\alpha-\gamma\alpha'', & \gamma\alpha'-\gamma'\alpha \\
\alpha'\beta''-\alpha''\beta', & \alpha''\beta-\alpha\beta'', & \alpha\beta'-\alpha'\beta
\end{array}
\]
manifesto in eandem formam transmutatur, in quam $f$ transit per hanc
\[
\begin{array}{ccc}
k, & 0, & 0 \\
0, & k, & 0 \\
0, & 0, & k
\end{array}
\]
sive in eam, quae oritur multiplicando singulos co\"efficientes formae $f$ per $kk$.
Hanc formam designabimus per $f'$.

IV. Prorsus simili modo probatur, formam $G$ per substitutionem $(S''')$
\[
\begin{array}{ccc}
\alpha, & \alpha', & \alpha'' \\
\beta, & \beta', & \beta'' \\
\gamma, & \gamma', & \gamma''
\end{array}
\]
\SourceLeaf{309}
transire in formam, quae oritur ex $F$, multiplicando singulos co\"efficientes per $kk$.
Hanc formam exprimemus per $F'$.

Substitutionem $S'''$ oriri dicemus \textit{per transpositionem} substitutionis $S$;
tunc manifesto $S$ rursus prodit ex transpositione substitutionis $S'''$; atque $S'$, $S''$
altera ex alterius transpositione. --- Substitutio $S'$ commode appellari potest
substitutioni $S$ \textit{adiuncta}, unde substitutioni $S'''$ adiuncta erit $S''$.

\Needspace{8\baselineskip}
\Article{269}
Si non modo forma $f$ implicat ipsam $g$, sed etiam haec illam, formae $f$, $g$
\textit{aequivalentes} vocabuntur. In hoc itaque casu non modo $D$ ipsum $E$ metietur,
sed etiam $E$ ipsum $D$, unde facile concluditur, esse debere $D=E$. Vice
versa autem, si forma $f$ implicat formam $g$ eiusdem determinantis, hae duae
formae erunt aequivalentes. Erit enim (adhibendo eadem signa ut in art.\ praec.
excipiendoque casum ubi $D=0$) $k=\pm1$, adeoque forma $f'$. in quam transit
$g$ per substitutionem $S''$, cum $f$ identica, sive $f$ sub $g$ contenta. Porro patet,
in hoc casu etiam formas $F$, $G$, ipsius $f$, $g$ adiunctas, inter se aequivalentes
fore, posterioremque in priorem transire per substitutionem $S'''$. Denique vice
versa, si formae $F$, $G$ aequivalentes esse \textit{supponuntur}, atque prior transit in
posteriorem per substitutionem $T$, etiam formae $f$, $g$ aequivalentes erunt, transibitque
$f$ in $g$ per substitutionem ipsi $T$ adiunctam, atque $g$ in $f$ per eam,
quae oritur ex transpositione substitutionis $T$. Nam per has duas substitutiones
resp.\ transit forma ipsi $F$ adiuncta in formam ipsi $G$ adiunctam atque haec in
illam; hae duae formae autem oriuntur ex $f$, $g$ multiplicando singulos co\"efficientes
per $D$; unde nullo negotio concluditur, per easdem substitutiones transire
$f$ in $g$. atque $g$ in $f$ resp.

\Needspace{8\baselineskip}
\Article{270}
Si forma ternaria $f$ formam ternariam $f'$ implicat, atque haec formam $f''$:
implicabit etiam $f$ ipsam $f''$. Facillime enim perspicitur, si transeat
\[
\begin{array}{c|c}
f\text{ in }f'\text{ per substitutionem} & f'\text{ in }f''\text{ per substitutionem} \\[2pt]
\begin{array}{ccc}
\alpha, & \beta, & \gamma \\
\alpha', & \beta', & \gamma' \\
\alpha'', & \beta'', & \gamma''
\end{array}
&
\begin{array}{ccc}
\delta, & \varepsilon, & \zeta \\
\delta', & \varepsilon', & \zeta' \\
\delta'', & \varepsilon'', & \zeta''
\end{array}
\end{array}
\]
\SourceLeaf{310}
$f$ transmutatum iri per substitutionem
\[
\begin{array}{ccc}
\alpha\delta+\beta\delta'+\gamma\delta'', & \alpha\varepsilon+\beta\varepsilon'+\gamma\varepsilon'', & \alpha\zeta+\beta\zeta'+\gamma\zeta'' \\
\alpha'\delta+\beta'\delta'+\gamma'\delta'', & \alpha'\varepsilon+\beta'\varepsilon'+\gamma'\varepsilon'', & \alpha'\zeta+\beta'\zeta'+\gamma'\zeta'' \\
\alpha''\delta+\beta''\delta'+\gamma''\delta'', & \alpha''\varepsilon+\beta''\varepsilon'+\gamma''\varepsilon'', & \alpha''\zeta+\beta''\zeta'+\gamma''\zeta''
\end{array}
\]

In eo itaque casu, ubi $f$ aequivalet ipsi $f'$, atque $f'$ ipsi $f''$, forma $f$ etiam
formae $f''$ aequivalebit. --- Ceterum sponte manifestum est, quomodo haec theoremata
ad plures formas sint applicanda.

\Needspace{8\baselineskip}
\Article{271}
Hinc iam patet, omnes formas ternarias, perinde ac binarias, in \textit{classes}
distribui posse, referendo ad classem eandem formas aequivalentes, non-aequivalentes
ad diversas. Formae itaque determinantium diversorum certo ad classes
diversas pertinebunt, et proin classes infinite multae formarum ternariarum dabuntur;
formae autem ternariae eiusdem determinantis modo minorem modo maiorem
classium numerum efficiunt; quod vero tamquam proprietas palmaris harum
formarum est considerandum, \textit{omnes formae eiusdem determinantis dati semper constituunt
classium multitudinem f\,i\,n\,i\,t\,a\,m.} Evolutioni uberiori huius gravissimi theorematis
praemittenda est explicatio sequentis differentiae essentialis, quae inter
formas ternarias obtinet.

Quaedam formae ternariae ita sunt comparatae, ut per ipsas sine discrimine
repraesentari possint numeri positivi et negativi, \textit{e.\ g.}\ forma $xx+yy-zz$, quamobrem
\textit{formae indefinitae} vocabuntur. Contra per alias numeri negativi repraesentari
nequeunt, sed (praeter cifram quae prodit, ponendo singulas indeterminatas
$=0$) positivi tantum, ut $xx+yy+zz$, quare \textit{formae positivae} dicentur; denique
per alias numeri positivi repraesentari nequeunt, ut $-xx-yy-zz$, unde
appellabuntur \textit{formae negativae}; formae positivae et negativae nomine communi
\textit{formae definitae} dicentur. Ecce iam criteria generalia, per quae haec formarum
indoles discerni poterit.

Multiplicando formam ternariam
\[
f=axx+a'x'x'+a''x''x''+2bx'x''+2b'xx''+2b''xx'
\]
\SourceLeaf{311}
determinantis $D$ per $a$, denotandoque co\"efficientes formae ipsi $f$ adiunctae
perinde ut in art.\ 267 per $A$, $A'$, $A''$, $B$, $B'$, $B''$, prodit
\[
(ax+b''x'+b'x'')^2-A''x'x'+2Bx'x''-A'x''x''=g
\]
multiplicando denuo per $A'$, provenit
\[
A'(ax+b''x'+b'x'')^2-(A'x''-Bx')^2+aDx'x'=h
\]
Hinc statim concluditur, si tum $A'$ tum $aD$ sint numeri negativi, omnes valores
ipsius $h$ esse negativos, unde manifesto per formam $f$ tales tantummodo numeri
repraesentari poterunt, quorum signum oppositum est signo ipsius $aA'$, \textit{i.\ e.}\ identicum
cum signo ipsius $a$, sive oppositum signo ipsius $D$. In hoc itaque casu $f$
erit forma definita, et quidem positiva vel negativa, prout $a$ est positivus vel
negativus, sive prout $D$ est negativus vel positivus.

Si vero vel vterque $aD$, $A'$ est positivus, vel alter positivus alter negativus
(neuter $=0$), facile perspicitur, $h$ per debitam quantitatum $x$, $x'$, $x''$ determinationem
valores tum positivos tum negativos nancisci posse. Quare in hoc casu
$f$ valores tum eodem signo affectos ut $aA'$ tum opposito obtinere poterit, eritque
adeo forma indefinita.

Pro eo casu, ubi $A'=0$, neque vero $a=0$, fit
\[
g=(ax+b''x'+b'x'')^2-x'(A''x'-2Bx'')
\]
Tribuendo ipsi $x'$ valorem arbitrarium (qui tamen non $=0$), accipiendoque $x''$
ita ut $\frac{A''x'}{2B}-x''$ signum idem obtineat ut $Bx'$ (quod fieri posse facile perspicitur,
quum $B$ nequeat esse $=0$, hinc enim foret $BB-A'A''=aD=0$, adeoque
etiam $D=0$, quem casum excludimus), erit $x'(A''x'-2Bx'')$ quantitas positiva,
unde facile patet, $x$ ita determinari posse, ut $g$ obtineat valorem negativum.
Manifesto hi valores etiam ita accipi poterunt, ut, si desideretur, omnes sint
integri. Denique patet, si ipsis $x'$, $x''$ valores quicunque tribuantur, ipsum $x$
tam magnum accipi posse, ut $g$ fiat positivus. Hinc concluditur, in hoc casu
formam $f$ esse indefinitam.
